% ═══════════════════════════════════════════════════════════════════
% EN Part 5: Certificate H, protocol V1–V9, summaries, reproducibility
% ═══════════════════════════════════════════════════════════════════
\part{Certificate H, Protocol V1--V9, Summaries}

\section{Certificate H: the $N=15/30$ stand in the full Gross
$\Ga$-normalization}\label{sec:certH}

Certificate~H completes the arch of the program: it implements the
full Gross $\Ga$-normalization on the $N=15/30$ stand and closes the
$\Ga$-remainder of certificate~G. Six theorems H1--H6 are accepted on
all points of the protocol.

\begin{theorem}[H1 --- the cyclotomic lattice]\label{th:H1}
The CM lattice of the stand: $\QQ(\zeta_{30})=\QQ(\zeta_{15})$ ---
one stand, two $\Ga$-levels. The Riemann form is integral (Ramanujan
sums); the polarization type in Smith normal form is
$(1,1,5,5,15,15,15,15)$ --- integer indices; the discriminant is
\[
  \disc \;=\; 3^4\cdot 5^6 \;=\; 1265625 \;=\; 1125^2
\]
--- a perfect square with $\sqrt{\disc}=1125$ integral. The CM types
$(1,2,4,7)$ for $N=15$ and $(1,7,11,13)$ for $N=30$ lie in one field.
\end{theorem}

\begin{proof}
The Riemann form: the pairing matrix of characters is built from the
Ramanujan sums $c_d(k)$ --- integers; integrality follows from the
integrality of the Ramanujan sums. The SNF is computed by the exact
integer algorithm; the diagonal $(1,1,5,5,15,15,15,15)$ gives the
product $5^2\cdot15^4=1265625$ --- the match with $3^4\cdot5^6$
verified by factorization. The field: $\phi(30)=\phi(15)=8$, and the
isomorphism $\QQ(\zeta_{30})\cong\QQ(\zeta_{15})$ is given by
$\zeta_{30}\mapsto-\zeta_{15}^{8}$.
\end{proof}

\begin{theorem}[H2 --- the full $\Ga$-normalization of the volume]
\label{th:H2}
The volume in the full $\Ga$-normalization:
\[
  \mathrm{vol}_h \;=\; (2\pi)^{32}\;
  \prod_{k=1}^{N-1}\Ga\!\Bigl(\frac{k}{N}\Bigr)^{-c_k}
  \;=\;1125
  \qquad (N=15\ \text{and}\ N=30),
\]
with exact exponents $c_k$ (for $N=15$: $c_k=4$ at
$k=1,2,4,5,7,8,10,11,13,14$ and $c_k=6$ at $k=3,6,9,12$); the
residual of the identity is $1.01\cdot10^{-119}$ at deep precision;
independent recovery of the exponents by PSLQ deflation at $N=15$
gives the exact integer vector; at $N=30$ the identity is certified
by the closed form (the sine count). Moreover $\det V^2=1265625/256$.
\end{theorem}

\begin{proof}
The volume formula is the Chowla--Selberg product over characters:
each character contributes a $\Ga$-monomial, the exponents collect
the conductor multiplicities $h_d$ (Appendix~\ref{app:census}).
Computation at $dps=120$ gives $1125$ with residual $10^{-119}$; the
match with $\sqrt{\disc}$ of Theorem~\ref{th:H1} is exact. At $N=15$
the PSLQ deflation recovered the exponent vector matching the
analytic count; at $N=30$ the analytic sine count
($\Ga(k/30)\Ga((30-k)/30)=\pi/\sin(\pi k/30)$) closes the identity
without PSLQ.
\end{proof}

\begin{theorem}[H3 --- Fermat periods as $\Ga$-monomials]\label{th:H3}
Every period of the Fermat curve is the $\Ga$-monomial
\[
  P_{a,b}(r,s) \;=\; \frac{1}{N}\,
  \zeta_N^{ra+sb}\,\frac{\Ga(\frac aN)\Ga(\frac bN)}
  {\Ga(\frac{a+b}N)}
  \qquad (91\ \text{and}\ 406\ \text{periods}),
\]
quadratures along the curve paths agree to $10^{-121}$; the
$\mu_N\times\mu_N$-equivariance is exact; level $30$ = values
$\Ga(k/30)$ with odd $k$ --- not reducible to $15$: a genuine
$\Ga$-rung of its own (the ``$\pi/15$ OR $\pi/30$'' of the directive
are distinct).
\end{theorem}

\begin{proof}
The closed form is Lemma~\ref{lem:closedform}; certificate~B (check
V2) cross-checks it against independent tanh-sinh on $69$ tests with
deviation $3.9\cdot10^{-36}$, and the deep record V9 at $dps=70$
gives $1.7\cdot10^{-71}$; the full quadratures of certificate~H along
the curve paths push the control to $10^{-121}$. Equivariance is
Corollary~\ref{cor:equivar}. Non-reducibility of level $30$: values
$\Ga(k/30)$ at odd $k$ (e.g.\ $\Ga(1/30)$) are not expressible
through $\Ga(k/15)$ and sines by the inventory of
Lemma~\ref{lem:relations} --- the argument $1/30$ reduces neither by
the recurrence nor by reflection to denominator $15$.
\end{proof}

\begin{theorem}[H4 --- the Gross normalization of the
``fifteenths'']\label{th:H4}
The moduli $\bigl|\omega(\tau_i)^2\bigr|$ are $\Ga$-monomials with
rational exponents (PSLQ deflation); the relative phase
\[
  \psi \;=\; \frac{(1+3\sqrt5)+i(\sqrt{15}-\sqrt3)}{8}
  \;=\; \frac{R}{|R|},
  \qquad R=\frac{(2\sqrt5-3)+(\sqrt5-2)\sqrt{-3}}{2},
\]
is recovered by LLL in the class field; the modulus $|R|=3-\sqrt5$
is a unit of $\QQ(\sqrt5)$; the $j$-pair of level $d=15$ agrees with
certificate~G.
\end{theorem}

\begin{proof}
Moduli: the Chowla--Selberg normalization (G4) gives
$\Ga$-monomials; the exponents are recovered by deflation
independently of the analytic count. Phase: LLL in
$\QQ(\sqrt{-3},\sqrt5)$ on the vector $(\re\psi,\im\psi)$ recovers
the exact coordinates $\frac{1+3\sqrt5}{8}$ and
$\frac{\sqrt{15}-\sqrt3}{8}$; the check $\psi=R/|R|$ via
$R\cdot\bar R=(3-\sqrt5)^2$ is exact.
\end{proof}

\begin{theorem}[H5--H6 --- the ladder and the denominators]
\label{th:H56}
\emph{(H5)} The levels $\pi/7$, $\pi/15$, $\pi/30$ are rungs of one
$\mu_N$-quantum ladder: the $\lambda_1$-pairings agree over the
weights $4\sin(\pi k/N)$. \emph{(H6)} The denominator bound of the
stand is
\[
  Q_{\text{stand}} \;=\; 2N\cdot\max_i s_i \;=\; 480
\]
($s_i$ --- the invariant factors of H1); the rationalization of
certificate~F on the stand closes without a scan: all target values
have denominators $\le480$, and the uniqueness of (F1) is guaranteed
at working precision.
\end{theorem}

\begin{proof}
(H5) The ladder: the phases $\pi/7$ (Klein, certificate~C), $\pi/15$
and $\pi/30$ (this stand) are successive levels of the construction
$\delta=\pi/N$; the pairing weights $4\sin(\pi k/N)$ agree with the
Ramanujan tables of H1. (H6) By Theorem~F2 the denominator bound is
governed by the determinant: $|\disc|=3^4 5^6$ divides $Q^{2g}$; the
direct count $2\cdot30\cdot8=480$ gives the bound with margin; the
scan of certificate~F was not needed --- all measurements of the
stand fall inside the uniqueness intervals.
\end{proof}

\begin{statusbox}[certificate H --- summary]
H1--H6 accepted ($6/6$): the cyclotomic lattice (Ramanujan form,
SNF type $(1,1,5,5,15,15,15,15)$, $\disc=1125^2$, field
$\QQ(\zeta_{30})=\QQ(\zeta_{15})$); the full $\Ga$-normalization of
the volume $\mathrm{vol}_h=(2\pi)^{32}\prod\Ga(k/N)^{-c_k}=1125$;
the Fermat periods as $\Ga$-monomials ($91/406$, quadratures
$10^{-121}$); the Gross normalization $d=15$ (the phase $\psi$, LLL
in the class field); the ladder $\pi/7,\pi/15,\pi/30$; the
denominator bound $Q=480$. Runtime $\sim$2 minutes.
\end{statusbox}

\section{The computational verification protocol V1--V9}
\label{sec:protocol}

The protocol (\texttt{laboratory.py}; results in \texttt{reports/})
is built so as to use neither PSLQ, nor L-values, nor numerical
checks of the Riemann relations:
\begin{enumerate}[leftmargin=2em, itemsep=0.15em]
\item all structural claims are verified by \emph{exact integer
  arithmetic} (censuses, ranks by construction, sums over
  conductors);
\item analytic identities (the closed form) are cross-checked by
  \emph{direct independent integration} via tanh-sinh over smooth
  pieces with no $\Ga$-functions at all (the substitution $t=u^N/2$
  on the halves of the interval removes endpoint singularities);
\item the ``completeness'' of the generating set is verified by the
  \emph{rank} of the explicitly built functional matrix --- a linear
  algebra check, not an integer-relation search;
\item working precision $dps=35$, control deep record $dps=70$.
\end{enumerate}

\begin{table}[t]
\centering
\caption{The checks V1--V9 and their results.}
\label{tab:vprotocol}
\small
\begin{tabularx}{\textwidth}{@{}l >{\raggedright\arraybackslash}X >{\raggedright\arraybackslash}p{4.3cm}@{}}
\toprule
Check & Result & Status \\
\midrule
V1: character census & $91=1+6+84$;
  $406=1+6+9+30+84+276$ & proved (App.~\ref{app:census}) + exact
  arithmetic \\
V2: closed form vs tanh-sinh & max rel.\ dev.\
  $3.9\cdot10^{-36}$ ($69$ tests) & computed ($dps$ $35$) \\
V3: phases $\arg P=\frac{2\pi(ra+sb)}{N}$ & max dev.\ $0.0$ &
  proved (Lemma~\ref{lem:closedform}) \\
V4: $\mu_N\times\mu_N$-equivariance & max rel.\ dev.\
  $3.9\cdot10^{-36}$ & proved (Cor.~\ref{cor:equivar}) + computed \\
V5: chain integrity $\int_\gamma dh=0$ & max abs.\
  $1.5\cdot10^{-36}$ & computed \\
V6: reflection ladder & max rel.\ dev.\ $7.0\cdot10^{-36}$ &
  proved (Lemma~\ref{lem:reflection}) + computed \\
V7: completeness & numerical rank $=g$: $91$ and $406$; cond.\
  $8.6$; $18.2$ & computed \\
V8: DFT orthogonality & diag.\ $4.0\cdot10^{-35}$; off-diag.\
  $8.2\cdot10^{-37}$ & computed \\
V9: deep record ($dps$ $70$) & rel.\ dev.\ $1.7\cdot10^{-71}$ &
  computed \\
\bottomrule
\end{tabularx}
\end{table}

Reproducibility: one command \texttt{python3 laboratory.py} runs the
whole protocol and reproduces every number of this monograph.
Dependencies: \texttt{mpmath}, \texttt{numpy}, \texttt{matplotlib};
fully deterministic --- no randomness, no seeds. Artifacts: report
files in \texttt{reports/} (JSON + text) and tiled $600$~dpi plots.

\section{Summary tables}\label{sec:summary}

\begin{table}[t]
\centering
\caption{The ladder of stands: integer triples and statuses.}
\label{tab:ladder-full}
\small
\begin{tabularx}{\textwidth}{@{}l l l l >{\raggedright\arraybackslash}X@{}}
\toprule
Stand & Triple & Genus & Key numbers & Status \\
\midrule
Torus & $(4,1,4\pi^2)$ & $1$ & $\delta=\pi/4$; holonomy $=i$;
  $4\pi^2$ & exact derivation \\
K3 & $(4,22,\cdot)$ & $-$ & $\rho=20$; $(1,19)$; $\det=64$;
  $48$ lines & exact computation \\
Klein & $(7,22,3.338)$ & $3$ & $j=-3375$; $\Delta=-343$;
  $\tau=\tfrac{1+\sqrt{-7}}2$ & certificates B, C \\
\textbf{Fermat 15} & $(15,\cdot,\cdot)$ & $91$ &
  $1+6+84=91$; radical $b_{Ch}(15)$ & closed system \\
\textbf{Fermat 30} & $(30,\cdot,\cdot)$ & $406$ &
  $1+6+9+30+84+276$; radical $b_{Ch}(30)$ & closed system \\
\bottomrule
\end{tabularx}
\end{table}

\begin{table}[t]
\centering
\caption{Program certificates: coverage and verdicts.}
\label{tab:certs}
\small
\begin{tabularx}{\textwidth}{@{}l >{\raggedright\arraybackslash}X l@{}}
\toprule
Certificate & Content & Verdict \\
\midrule
A & explicit divisor with equations; exact certificate over $\QQ$;
  kernel $29$ & accepted \\
B & closing the kernel $29$ by periods; Theorem B1; blind spot
  $29+2\to0$; $51$ measurements on the $22$-dim $H^2$ & accepted \\
C & $\mu_4$-equivariance; $22=1+7+7+7$;
  $\mathrm{charpoly}(\lambda_1)=x^4-1$; the $\sqrt{-7}$ heptad &
  accepted \\
D & universal $\mu_4$ theorem: $[L,P]=0$; functoriality;
  semi-invariance over $\ZZ[i]$ & accepted \\
E & flow termination: E1--E4; $t^*=\lcm(\dots)$; cross-check
  $48/212/432/114$ & accepted ($8/8$) \\
F & rationalization: F1--F4; a priori bound $Q=256$; LLL:
  $4X^2-7$; $j=-3375$ & accepted ($5/5$) \\
G & Hodge lattice indices: G1--G6; Chowla--Selberg; $\pi/15$,
  $\pi/30$ & accepted ($6/6$) \\
H & the $N=15/30$ stand: H1--H6; SNF $(1,1,5,5,15,15,15,15)$;
  $\disc=1125^2$; $\mathrm{vol}_h=1125$; $Q=480$ &
  accepted ($6/6$) \\
\bottomrule
\end{tabularx}
\end{table}

In total the program contains: $3$ calibration stands (torus, K3,
Klein), $2$ cyclotomic levels ($N=15$, $N=30$) with a closed system
of $8$ lemmas and $6$ theorems, $8$ certificates with full protocols,
$9$ checks V1--V9, and $2$ errata records (E8 $28/36$, the E6 case).
All results are reproducible by the laboratory \texttt{laboratory.py}
in one command and are verified by the multilingual panel (Python,
Julia, Fortran, C, Rust) and by the Lean kernel (exact integer
identities).

\section{Reproducibility guide}\label{sec:labmap}

Every numerical result of the monograph has a direct source in the
\texttt{laboratory.py} menu. The mapping:

\begin{center}
\small
\begin{tabularx}{\textwidth}{@{}l l >{\raggedright\arraybackslash}X@{}}
\toprule
Section & Numbers & Laboratory item \\
\midrule
\ref{sec:torus} (torus) & $(4,1,4\pi^2)$; $\gamma$; $\Delta$ &
``Stands $\to$ Torus'' \\
\ref{sec:k3} (K3) & $\rho=20$; $(1,19)$; $64$ & ``Stands $\to$ K3'' \\
\ref{sec:e8} (errata) & $28/36$; $64$ & ``Stands $\to$ Errata E8'' \\
\ref{sec:binary} (code) & $48/212/432/114$ & ``Stands $\to$ Binary
code'' \\
\ref{sec:certA}--\ref{sec:certB} & kernel $29$; assembly $[Z]$ &
``Certificates $\to$ A, B'' \\
\ref{sec:certC}--\ref{sec:certD} & $22=1+7+7+7$; $[L,P]=0$ &
``Certificates $\to$ C, D'' \\
\ref{sec:certE} & $t^*=48$ & ``Certificates $\to$ E'' \\
\ref{sec:certF} & $Q=256$; $4X^2-7$ & ``Certificates $\to$ F'' \\
\ref{sec:certG} & Chowla--Selberg; $\pi/15$, $\pi/30$ &
``Certificates $\to$ G'' \\
\ref{sec:klein} (Klein) & $j=-3375$; $\Omega$ & ``Stands $\to$
Klein'' \\
\ref{sec:cycles}--\ref{sec:indices} & V1--V9; SNF; $\mathrm{vol}_h$
& ``Stand N=15/30'' (full) \\
\ref{sec:certH} (H) & H1--H6 & ``Certificates $\to$ H'' \\
\bottomrule
\end{tabularx}
\end{center}

\subsection{Quick start}

\begin{enumerate}[leftmargin=2em, itemsep=0.15em]
\item Install dependencies: \texttt{pip install mpmath numpy
  matplotlib}.
\item Run the lab: \texttt{python3 laboratory.py}.
\item Choose the interface language (Russian/English).
\item ``Full protocol V1--V9'' --- reproduces
  Table~\ref{tab:vprotocol} ($\sim$3 minutes).
\item ``Reports and plots'' --- creates \texttt{reports/}: JSON
  reports, a text log, and $600$~dpi tiled plots.
\item ``Multilingual verification'' --- if \texttt{julia},
  \texttt{gfortran/clang}, \texttt{cargo}, \texttt{lean} are
  present, runs the parallel implementations and compares verdicts.
\end{enumerate}

\subsection{Comparison rules}

When comparing monograph numbers with lab output, the threshold rule
applies: a value is reproduced if the relative deviation is within
the declared protocol threshold ($10^{-30}$ at $dps=35$,
$10^{-65}$ at $dps=70$; integers --- exact equality). Thresholds are
embedded in the reports; the lab prints PASS/FAIL automatically, and
the summary verdict ``all checks PASS'' is equivalent to reproducing
the section.

\section{Verification in five languages and the Lean kernel}
\label{sec:multilang}

\subsection{Roles of the languages}

\begin{center}
\small
\begin{tabularx}{\textwidth}{@{}l l >{\raggedright\arraybackslash}X@{}}
\toprule
Language & File & Role \\
\midrule
Python & \texttt{laboratory.py} & the full laboratory: protocols,
reports, plots \\
Julia & \texttt{verification/julia/verify\_hodge.jl} &
high-precision periods and phases \\
Fortran & \texttt{verification/fortran/verify\_hodge.f90} &
census and rank arithmetic \\
C & \texttt{verification/c/verify\_hodge.c} & fast integer checks
(Bareiss) \\
Rust & \texttt{verification/rust/verify\_hodge.rs} & safe integer
arithmetic \\
Lean~4 & \texttt{verification/lean/HodgeLaboratory.lean} & exact
identities in the kernel \\
\bottomrule
\end{tabularx}
\end{center}

Each implementation checks the same base set of identities: the
genera $91$ and $406$; the censuses $1+6+84$ and
$1+6+9+30+84+276$; the discriminant $1125^2$; the SNF product; the
$b_{Ch}$ radicals; the termination $t^*=\lcm$; the cross-checks
$48/212/432/114$; the Klein $j=-3375$. The agreement of five
independent implementations and the Lean kernel is the strongest
reproducibility certificate: an error repeated simultaneously in
five ecosystems is practically excluded.

\subsection{Lean: what exactly the kernel proves}

The file \texttt{HodgeLaboratory.lean} contains formalized
propositions checked by the Lean~4 kernel without axioms beyond the
standard ones:
\begin{itemize}[leftmargin=1.6em, itemsep=0.1em]
\item the Fermat genus: $(15-1)*(15-2)/2=91$,
  $(30-1)*(30-2)/2=406$;
\item the censuses by conductors: exact M\"obius sums
  $1+6+84=91$, $1+6+9+30+84+276=406$;
\item the stand discriminant $3^4*5^6=1265625=1125^2$;
\item the SNF product $(1,1,5,5,15,15,15,15)$ equals $1265625$;
\item the radicals $b_{Ch}(15)$ and $b_{Ch}(30)$ as expressions in
  $\sqrt5$, $\sqrt3$, $\sqrt{10\pm2\sqrt5}$;
\item termination: $t^*=\lcm(W/g, H/g)$ and the case $t^*=48$ for
  the $48\times48$ grid;
\item the binary-code cross-check: $48,212,432,114$;
\item the Klein integer invariants: $105^3=343*3375$,
  $j=-3375=-15^3$.
\end{itemize}
The kernel verifies these propositions by complete computation
(\texttt{native\_decide}, \texttt{norm\_num}, \texttt{ring}) ---
machine verification in the strictest sense; the run is the menu
item ``Lean verification'' or the command
\texttt{lean HodgeLaboratory.lean}.

\subsection{Optional environments}

All five extra environments are optional: the laboratory is fully
functional on Python alone; with the environments present it
automatically extends the verdicts. Installation in brief:

\begin{itemize}[leftmargin=1.6em, itemsep=0.1em]
\item Julia: \texttt{juliaup} or the distribution package manager;
  run: \texttt{julia verification/julia/verify\_hodge.jl}.
\item Fortran: \texttt{gfortran -O2 -o verify\_hodge
  verify\_hodge.f90 \&\& ./verify\_hodge}.
\item C: \texttt{gcc -O2 -o verify\_hodge verify\_hodge.c -lm
  \&\& ./verify\_hodge}.
\item Rust: \texttt{rustc -O verification/rust/verify\_hodge.rs -o
  verify\_hodge \&\& ./verify\_hodge} (or \texttt{cargo run}).
\item Lean: \texttt{elan} (toolchain \texttt{leanprover/lean4});
  run: \texttt{lean verification/lean/HodgeLaboratory.lean}.
\end{itemize}

\subsection{The full reproduction checklist}

\begin{enumerate}[leftmargin=2em, itemsep=0.15em]
\item \texttt{python3 laboratory.py} $\to$ ``Full protocol'' $\to$
  summary PASS.
\item Reports in \texttt{reports/}: JSON files match the references
  in \texttt{results/} by keys and thresholds.
\item Plots \texttt{reports/plots/}: eight $600$~dpi tiles generated
  without errors.
\item (Optional) five languages: all verdicts PASS, counters match.
\item (Optional) Lean: the file compiles without errors, all
  \texttt{example}s accepted by the kernel.
\item GitHub Actions (after push): \texttt{ci.yml} green --- the
  Python protocol run on three Python versions.
\end{enumerate}

After passing the checklist the monograph is fully reproduced: every
number of Tables~\ref{tab:vprotocol}--\ref{tab:certs} and of all
sections is confirmed by an independent run.
